% Root and geometric_routes independently approved the complete proof.
% Integrated as FF.50; retained separately as the reviewed source.
\subsection{Integral norms and a higher-genus slice}

\begin{claim}{FF.50}{Proved: Kummer rigidity for an explicit genus-two slice}
\label{clm:FFhighergenusnorm}
Retain the pointwise pinned automorphism hypotheses of
Claims~\ref{clm:FFellipticshear} and the preceding
genus-zero and Lattes results, over a finite extension
\(E/B_p\). Let \(C\) be the smooth projective curve
\[
 C:\quad y^2=(u^2-1)(u^2-2)(u^2-3),\qquad T=u^2.
                                                        \tag{FF.50.1}
\]
This is a genus-two Galois cover of \(\mathbb P^1_T\)
with group \(V_4\). Its transported actual point
function, including every completed point, is realized
by a unique scheme isomorphism over \(\mathbb P^1_T\).
Using that isomorphism to identify its generic action,
the automorphism fixes
\[
                    H^1(G_{L(C)},\mu_m)
                                                        \tag{FF.50.2}
\]
for every finite extension \(L/E\) and every positive
integer \(m\).

More generally, the same Kummer conclusion holds for a
realized smooth proper slice \(C\) whenever it has
coherently realized connected finite etale covers
\(D_s\to C\), with completed-point compatibility,
whose pushforward images in integral geometric first
homology are contained in
\(\ell^sH_1(C,\mathbb Z_\ell)\), for arbitrarily
large \(s\) and each prime \(\ell\).
\end{claim}

The completed-point hypothesis in the general criterion
is substantive. Preserving the isomorphism class of an
Abel cover does not establish it. The specific curve
(FF.50.1) meets the hypothesis by its two elliptic
quotients, as proved below. This is a positive-genus
Kummer constraint; it does not identify all further
nonabelian covers or prove reconstruction over \(B_p\).

\begin{proof}
\emph{1. The integral norm criterion.}
Work first over an algebraic closure of the constants.
For every coherently realized curve \(X\), use its
actual-point scheme isomorphism to identify the
transported generic group with \(G_{\overline E(X)}\).
The induced action fixes each geometric place and its
inertia exponent, the latter by restriction of the
base exponent one. Hence for each rational function
\(f\ne0\) its Kummer error is unramified:
\[
 b_X(f):=\alpha_X\kappa_\ell(f)-\kappa_\ell(f)
             \in H^1(X_{\overline E},\mathbb Z_\ell(1)).
                                                        \tag{FF.50.3}
\]
The localization sequence proves this integrally:
the residues at all places agree, so their difference
lies in proper first cohomology.

Let \(\rho:D\to C\) be any connected finite etale
cover in the stated coherent family. The norm
\(\overline E(D)^\times\to\overline E(C)^\times\)
is surjective. Here is the dimension proof needed for
that assertion. The function field of a curve over an
algebraically closed field is \(C_1\): for a
homogeneous form of degree \(d\) in \(N>d\) variables,
choose the variables in \(H^0(C,\mathcal O_C(mA))\),
where \(A\) is a fixed positive divisor. After clearing
the fixed coefficient poles, the equation is a system
of at most \(dm\deg A+O(1)\) homogeneous equations in
\(Nm\deg A+O(1)\) scalar variables, by Riemann--Roch.
For sufficiently large \(m\) there are fewer equations
than variables, so the projective dimension theorem
gives a nonzero common zero over the algebraically
closed constant field. Apply this to
\[
               N_{\overline E(D)/\overline E(C)}(v)-f z^d,
               \qquad d=[\overline E(D):\overline E(C)].
                                                        \tag{FF.50.4}
\]
It has degree \(d\) in \(d+1\) variables. A nonzero
solution has \(z\ne0\), since the norm of a nonzero
field element is nonzero. Dividing by \(z\) supplies
\(h\in\overline E(D)^\times\) with \(N(h)=f\).
All these data descend to a finite constant extension.
This is the usual Tsen argument, included here in the
precise form used.

Functoriality of corestriction and of the coherent
scheme realizations gives
\[
                   b_C(f)=\operatorname{Tr}_\rho b_D(h).
                                                        \tag{FF.50.5}
\]
The transfer direction matters: under integral
Poincare duality for the two compact curves,
\(\operatorname{Tr}_\rho:H^1(D,\mathbb Z_\ell(1))
\to H^1(C,\mathbb Z_\ell(1))\) corresponds to
\(\rho_*:H_1(D,\mathbb Z_\ell)\to
H_1(C,\mathbb Z_\ell)\). Thus the hypothesis on
the image of homology makes (FF.50.5) divisible by
\(\ell^s\) in the fixed free \(\mathbb Z_\ell\)-module
\(H^1(C,\mathbb Z_\ell(1))\). Varying \(s\) forces
\(b_C(f)=0\).

For the full Abel mod-\(n\) cover of \(C\), the
homology image is exactly \(nH_1(C,\mathbb Z)\):
its fundamental subgroup is the kernel of
\(\pi_1(C)\to H_1(C,\mathbb Z/n)\).
Every element of its homology image is divisible by
\(n\), and the \(n\)-th power of a loop realizing
any integral homology basis element gives the reverse
inclusion. Therefore these Abel covers would satisfy
the transfer requirement if their completed-point
compatibility were known. No such compatibility for
arbitrary curves is being assumed.

\emph{2. The three quotients of the specified curve.}
The six roots of the polynomial in (FF.50.1) are
distinct, so its double cover of \(\mathbb P^1_u\)
has genus two. The independent quadratic equations
\(u^2=T\) and \(y^2=(T-1)(T-2)(T-3)\) give
the \(V_4\)-cover over \(T\). Its three quadratic
quotients are
\[
 \begin{gathered}
 \mathbb P^1_u,\qquad A_1:\ y^2=(T-1)(T-2)(T-3),\\
 A_2:\ z^2=T(T-1)(T-2)(T-3),\qquad z=uy.
 \end{gathered}
                                                        \tag{FF.50.6}
\]
The latter two are elliptic curves, each with origin
at its branch point \(T=1\). Their maps to
\(\mathbb P^1_T\) are the quotients by \(\pm1\).
The point \(P_0=(u,y)=(1,0)\) on \(C\) maps to
both origins.

The previous genus-zero and Lattes results realize
all three quotient objects naturally, including
their completed-point maps. The three maps from
\(C\) jointly separate every geometric completed
point. For a group-theoretic verification, a
geometric point has cyclic inertia \(D\leq V_4\),
and equality of its images in a quotient by one
of the three order-two subgroups \(H_i\) means
that the relative translating element lies in
\(H_iD\). One has
\(\bigcap_iH_iD=D\): this is immediate if
\(D=1\), and if \(|D|=2\) one of the \(H_i\)
is exactly \(D\). Thus no distinct completed
points have all three images equal. Naturality
on the branch-complement opens, the unique
extension of a curve isomorphism to its smooth
completion, and this separation show that the
realization of \(C\) agrees with its entire
completed-point bijection. This step does not
assume that arbitrary normalized products
separate boundary branches.

\emph{3. Elliptic division covers provide the required
homology divisibility.}
Let \(\rho_i:C\to A_i\) be the two elliptic quotient
maps, based at \(P_0\). They induce homomorphisms
\[
 \varphi:J_C\longrightarrow A_1\times A_2,
 \qquad
 \psi:A_1\times A_2\longrightarrow J_C,
                                                        \tag{FF.50.7}
\]
where \(\varphi\) is the product of the pushforwards
and \(\psi\) the sum of the pullbacks. Their composite
is
\[
                              \psi\varphi=[2]_{J_C}.
                                                        \tag{FF.50.8}
\]
Indeed if \(\tau(u,y)=(-u,y)\) and
\(\iota(u,y)=(u,-y)\), the two degree-two norm
identities give the endomorphisms \(1+\tau\) and
\(1+\iota\tau\) on \(J_C\). The hyperelliptic
involution \(\iota\) acts as \(-1\), and their
sum is \(2\). In particular \(\varphi\) has
finite kernel killed by two and is an isogeny,
since both dimensions are two.

For \(n\geq1\), take the connected component
containing \((P_0,0,0)\) of
\[
 D_n=C\mathbin{\times_{A_1\times A_2,[n]}}
                                  (A_1\times A_2).
                                                        \tag{FF.50.9}
\]
It is a finite etale cover of \(C\). Its two elliptic
factors, considered over \(\mathbb P^1_T\), are
the Galois genus-one Lattes covers
\(A_i\xrightarrow{[n]}A_i\to\mathbb P^1_T\),
after a finite constant extension making the torsion
rational. They and all their displayed arrows have
the coherent completed-point realizations already
proved. Ordinary fiber products and the selected
component now give that same compatibility for
\(D_n\). This time there is no normalization
ambiguity: (FF.50.9) is etale over the smooth curve
\(C\), hence smooth, and its geometric points are
exactly the indicated tuples. The projection to
\(C\) and the two elliptic coordinates therefore
separate all its completed points.

Let \(\Lambda=H_1(C(\mathbb C),\mathbb Z)\) and
\(\Lambda_A=H_1((A_1\times A_2)(\mathbb C),\mathbb Z)\),
using any embedding of the finite coefficient field
into \(\mathbb C\). The covering subgroup in
(FF.50.9) consists exactly of the loops whose
\(\varphi_*\)-image lies in \(n\Lambda_A\).
Consequently each element \(v\) in the image of
\(H_1(D_n,\mathbb Z)\to\Lambda\) satisfies
\[
                 \varphi_*v\in n\Lambda_A,
                 \qquad 2v=\psi_*\varphi_*v\in n\Lambda.
                                                        \tag{FF.50.10}
\]
For \(n=\ell^s\), this puts the image in
\(\ell^s\Lambda\) if \(\ell\ne2\), and in
\(2^{s-1}\Lambda\) if \(\ell=2\). Integral
comparison gives the same assertion for etale
\(\mathbb Z_\ell\)-homology. The loss of one power
at two is uniform and harmless. Step~1 therefore
proves zero geometric Kummer error on \(C\) at
every prime, including two.

\emph{4. Arithmetic constants and finite coefficients.}
Fix a finite \(L/E\) and \(f\in L(C)^\times\).
Since \(\mu_\infty\subset L\), its compatible
\(\ell\)-adic Kummer class is a genuine continuous
character. The difference under the identified
generic automorphism is zero on the geometric
kernel by the preceding steps, hence is the
inflation of a character
\(\gamma:G_L\to\mathbb Z_\ell(1)\).
Corestriction to \(L(T)\) is zero, since rational
Kummer classes are fixed. On the inflated character
it is multiplication by \([L(C):L(T)]=4\).
The group \(\operatorname{Hom}_{\mathrm{cont}}
(G_L,\mathbb Z_\ell(1))\) is torsion-free, so
\(\gamma=0\), including when \(\ell=2\).
Reducing the integral Kummer character gives every
finite Kummer class, and combining the primes
proves (FF.50.2).
For the general criterion the identical argument uses
the positive degree \([L(C):L(T)]\) of its finite
slice in place of four; no divisibility assumption
on that degree is needed.
\end{proof}

\paragraph{Exact limitation and the next obstruction.}
The construction works because the Jacobian of this
curve admits an isogeny to elliptic quotients which
are themselves realized covers of the same rational
base. More generally the argument applies whenever
finitely many such quotient maps induce an injection
on rational first homology and their division
fiber products have completed-point compatibility;
a fixed isogeny denominator only changes the bounded
loss in (FF.50.10). For a Jacobian with an absolutely
simple factor of dimension greater than one, these
elliptic maps do not provide that injection. The
obstruction cannot be removed merely by taking a
finite cover \(D\to C\): norm composed with
pullback is multiplication by its degree on
\(J_C\), so \(J_C\) remains an isogeny factor
of \(J_D\). If the latter were built entirely
from elliptic factors, the former would be as well.
The
fixed-prime configuration theorem establishes a
completion topology, not the required actual point
maps on these higher-dimensional Abel covers.
Thus the remaining compatibility problem has not
been solved for an arbitrary positive-genus slice.
